\section*{Reproducing the results}

\texttt{pip install -r requirements.txt} installs the versions used here. \texttt{python study/run\_study.py} runs the capacity study and takes about five minutes on one processor core, and \texttt{python study/analyze.py} writes the summaries to \texttt{study/results/} and the figures to \texttt{study/figures/}. Every table and nearly every result number in this paper from the corrected baseline and the capacity study is generated from those files, from the committed baseline run in \texttt{hopfield\_results/}, and from the data files for the statistics of the characters, by \texttt{paper/make\_numbers.py}. \texttt{python scripts/kamp\_hw2.py} reruns the assignment's baseline and the dataset statistics on the files in \texttt{data/}, into new folders; move the new baseline run out of \texttt{hopfield\_results/} before running \texttt{make\_numbers.py}, which stops when it finds two. The original submission's numbers, in the revision note and in Section~\ref{sec:characters}, are kept as history: they come from its report and result files, which are not part of this repository. A few constants and counts are written into the text, such as the training density in the overlap formula of Section~\ref{sec:data} and the 18 of 18 in Section~\ref{sec:characters}.
